% Fichier complété par tools/complete_missing_corrections.py.
% Source : CIN/CIN-02-VitesseAcceleration/13_TransfoMouvement/13_TransfoMouvement.tex
% Statut : rédigé (2 question(s)).
\begin{corrigebox}[Corrigé rédigé]
\CorrigeQuestion{1}
Le solide 3 est en translation suivant son axe :
                $\vec\Omega_{3/0}=\vec0$ et $\vec V(B,3/0)=\dot\lambda\,\vec i_2$.
                Donc $\{\mathcal V(3/0)\}_B=\left\{\begin{array}{c}\vec0\\
                \dot\lambda\vec i_2\end{array}\right\}_B$. La valeur de $\dot\lambda$
                s'obtient en dérivant la fermeture géométrique.
\CorrigeQuestion{2}
Comme la direction de la glissière est fixe,
                $\vec\Gamma(B,3/0)=\ddot\lambda\,\vec i_2$. Si l'axe $\vec i_2$ est
                mobile, il faut ajouter le terme $\dot\lambda\,\vec\Omega_{2/0}\wedge
                \vec i_2$.
\end{corrigebox}
